% section: おまけ：定数数列

\section{おまけ：定数数列}\label{節：おまけ：定数数列}
  \bookterm{sequence}{数列}の種類には定数\bookterm{sequence}{数列}というものがある.
  時折,\bookterm{recurrence}{漸化式}を解く中で定数\bookterm{sequence}{数列}が生まれることがある.
  先程の解説から引用しよう.
  \begin{align*}
    &\frac{a_{n}}{n(n+1)}=\frac{a_{n-1}}{(n-1)(n)}=\frac{a_{n-2}}{(n-2)(n-1)}\\
    &\qquad=\cdots=\frac{a_{1}}{1(1+1)}
  \end{align*}

  これはまさに定数\bookterm{sequence}{数列}の良い例である.
  少し式が複雑なので,
  \[
    b_n=\frac{a_{n}}{n(n+1)}
  \]
  とおくと式は
  \[
    b_n=b_{n-1}=b_{n-2}=\cdots=b_1
  \]
  となる.
  ずっと同じ値を取るのだ.
  当たり前過ぎて何を言っているんだ,と思うかもしれないが急に出てきてパニックにならないようにここで念のため触れておく.
  この「一定の値を取り続ける\bookterm{sequence}{数列}」という発想は,\sectionref{節：特性方程式型}で登場する\bookterm{characteristic}{特性方程式}の解$\alpha$(\bookterm{recurrence}{漸化式}$a_{n+1}=pa_n+q$自身を定数\bookterm{sequence}{数列}にしてしまう値)にもつながっている.
